% ==================== optimality ====================
\section{Условия оптимальности}
\label{sec:optimality}

В предыдущем разделе получены линейные уравнения перехода. Теперь
необходимо выбрать число импульсов, их положения и векторы так, чтобы
минимизировать суммарную характеристическую скорость. Сначала выведем
условия оптимальности непосредственно для этой задачи. Затем установим
их связь с сопряжёнными уравнениями Лоудена для точного движения в
центральном поле \cite[гл.~3, \S~5.1]{lawden}.

\subsection{Задача с произвольным числом импульсов}

В безразмерных переменных требуется найти
\begin {equation}
 \begin {gathered}
 J_* = \min_{N,\,\varphi_j,\,\Delta\widehat{\boldsymbol V}_j^{\text{\normalfont имп}}}
       \sum_{j=1}^{N}\lVert\Delta\widehat{\boldsymbol V}_j^{\text{\normalfont имп}}\rVert,\\
 \sum_{j=1}^{N}B(\varphi_j)\Delta\widehat{\boldsymbol V}_j^{\text{\normalfont имп}}=\boldsymbol b,
 \qquad \boldsymbol b=({\Delta a},{\Delta e_x},{\Delta e_y},0,{\Delta i})^{\mathsf T}.
 \end {gathered}
 \label {eq:linear-optimization}
\end {equation}
Здесь $B(\varphi)$ --- матрица из уравнений перехода,
$\varphi_j$ определены по модулю $2\pi$,
$\Delta\widehat{\boldsymbol V}_j^{\text{\normalfont имп}}\in\R^3$, а число $N$ заранее не задано.
Нулевой импульс можно удалить, поэтому при обсуждении отдельного
манёвра будем считать ${\Delta\widehat V_j^{\text{\normalfont имп}}}=\lVert\Delta\widehat{\boldsymbol V}_j^{\text{\normalfont имп}}\rVert>0$.
При $\boldsymbol b=0$ решение очевидно: $J_*=0$ и импульсы отсутствуют.
Далее рассматривается $\boldsymbol b\ne0$.

Свобода выбора фаз в \eqref{eq:linear-optimization} связана с принятой
постановкой: время перехода и конечная фаза не ограничены. В линейной
модели импульсы с заданными фазами можно расположить в хронологическом
порядке, при необходимости добавляя обороты. Этот факт не означает,
что приближение остаётся равномерно точным при произвольно большом
времени движения. Погрешность восстановления реальной траектории
оценивается отдельно.

Представим ненулевой импульс в виде
\begin {equation}
 \Delta\widehat{\boldsymbol V}_j^{\text{\normalfont имп}}={\Delta\widehat V_j^{\text{\normalfont имп}}}\boldsymbol d_j,
 \qquad {\Delta\widehat V_j^{\text{\normalfont имп}}}>0,\qquad \lVert\boldsymbol d_j\rVert=1.
\end {equation}
Вектор $B(\varphi_j)\boldsymbol d_j\in\R^5$ показывает, какие изменения
параметров орбиты даёт единичная затрата характеристической скорости
в выбранной точке и направлении. Введём множество таких изменений
\begin {equation}
 \mathcal S=
 \bigl\{B(\varphi)\boldsymbol d:
          0\leq\varphi\leq2\pi,\ \lVert\boldsymbol d\rVert=1\bigr\}
\end {equation}
и его выпуклую оболочку $\mathcal C=\operatorname {conv}\mathcal S$.
Под выпуклой оболочкой понимается множество всех конечных средних
$\sum_j w_j\boldsymbol s_j$, где
$\boldsymbol s_j\in\mathcal S$, $w_j\geq0$ и $\sum_jw_j=1$.
Если суммарная затрата равна $J$, то
\begin {equation}
 \frac{{\boldsymbol b}}{{J}}
   =\sum_j\frac{{{\Delta\widehat V_j^{\text{\normalfont имп}}}}}{{J}}B(\varphi_j)\boldsymbol d_j
   \in\mathcal C.
 \label {eq:reachable-convex}
\end {equation}
Обратно, любое такое среднее задаёт набор импульсов
$\Delta\widehat{\boldsymbol V}_j^{\text{\normalfont имп}}=Jw_j\boldsymbol d_j$ с затратой $J$.
Следовательно, поиск оптимального управления эквивалентен поиску
наименьшего множителя $J$, для которого $\boldsymbol b\in J\mathcal C$.

Множество $\mathcal C$ компактно, симметрично относительно начала
координат и содержит начало координат внутри. Последнее свойство
существенно: все пять требуемых приращений доступны при конечной
затрате. Его можно проверить без анализа отдельных манёвров.
Если вектор $\boldsymbol l\in\R^5$ ортогонален всем элементам
$\mathcal S$, то $B(\varphi)^{\mathsf T}\boldsymbol l=0$ при любом $\varphi$.
Из первой и второй компонент этого тождества следуют
$l_x=l_y=l_a=0$, а из третьей --- $l_s=l_c=0$.
Таким образом, $\mathcal S$ порождает всё пространство $\R^5$;
его симметричная выпуклая оболочка имеет непустую внутренность.
Компактность можно пояснить следующим образом. Любое выпуклое
сочетание более чем шести точек в $\R^5$ допускает сокращение числа
слагаемых за счёт их аффинной зависимости. Повторяя сокращение,
получаем не более шести точек. Множество всех таких сочетаний
компактно, поскольку $\mathcal S$ компактно.
Число шесть здесь является общей геометрической оценкой; специальная
структура орбитального перехода позволяет уменьшить его.

Поэтому минимум $J_*>0$ существует, а точка
$\boldsymbol b/J_*$ лежит на границе $\mathcal C$.
Через граничную точку выпуклого тела можно провести опорную
гиперплоскость: всё тело находится по одну сторону от неё.
В рассматриваемом случае это означает существование
$\boldsymbol l\ne0$, для которого после выбора масштаба
\begin {equation}
 \boldsymbol l^{\mathsf T}\boldsymbol s\leq1
       \quad(\boldsymbol s\in\mathcal C),
 \qquad
 \boldsymbol l^{\mathsf T}\boldsymbol b=J_*.
 \label {eq:supporting-plane}
\end {equation}
Положительность нормирующей константы следует из того, что начало
координат находится внутри $\mathcal C$.
Геометрический факт об опорной плоскости можно получить, приближая
граничную точку внешними точками. Для каждой внешней точки выбирают
ближайшую точку выпуклого тела. Вектор между ними образует неострый
угол с любым направлением из ближайшей точки внутрь тела; иначе
расстояние можно было бы уменьшить. Предельное положение
перпендикулярной ему плоскости и даёт опорную плоскость.

\subsection{Базис-вектор и необходимое условие для всех допустимых фаз}

Множители $l_a,l_x,l_y,l_s,l_c$ задают линейную оценку изменений
параметров орбиты. Если единичное направление импульса равно
$\boldsymbol d$, импульс величины $\varepsilon$ создаёт изменение
$\delta\boldsymbol b=\varepsilon B(\varphi)\boldsymbol d$.
Его оценка равна
\[
 \boldsymbol l^{\mathsf T}\delta\boldsymbol b
 =\varepsilon\bigl(B(\varphi)^{\mathsf T}\boldsymbol l\bigr)^{\mathsf T}
 \boldsymbol d.
\]
Таким образом, множители граничных условий можно перенести из
пространства орбитальных параметров в пространство трёх компонент
импульса. Полученный вектор и задаёт оценку эффективности его направления.

Определим \emph{базис-вектор}
\begin {equation}
 \boldsymbol p(\varphi)=B(\varphi)^{\mathsf T}\boldsymbol l,
 \qquad
 \boldsymbol l=(l_a,l_x,l_y,l_s,l_c)^{\mathsf T}.
 \label {eq:primer-definition}
\end {equation}
В принятой нормировке все его компоненты безразмерны:
\begin {equation}
 \boxed{\begin {aligned}
 p_r(\varphi)&=l_x\sin\varphi-l_y\cos\varphi,\\
 p_t(\varphi)&=2l_a+2l_x\cos\varphi+2l_y\sin\varphi,\\
 p_z(\varphi)&=l_s\sin\varphi+l_c\cos\varphi.
 \end {aligned}}
 \label {eq:primer-components}
\end {equation}
Пять коэффициентов $\boldsymbol l$ постоянны. Это связано с тем, что
в выбранных орбитальных переменных линейная матрица действия импульса
зависит от фазы, но не от уже накопленных приращений орбиты.

Первое неравенство \eqref{eq:supporting-plane}, применённое к
$\boldsymbol s=B(\varphi)\boldsymbol d$, принимает вид
\begin {equation}
 \boldsymbol p(\varphi)^{\mathsf T}\boldsymbol d\leq1
        \qquad\text {при любом }\lVert\boldsymbol d\rVert=1.
\end {equation}
Наибольшее скалярное произведение с единичным вектором равно модулю
самого вектора. Следовательно,
\begin {equation}
 \lVert\boldsymbol p(\varphi)\rVert\leq1
       \qquad\text {при всех }\varphi.
 \label {eq:primer-global-bound}
\end {equation}
В оптимальном разложении \eqref{eq:reachable-convex} среднее
скалярных произведений равно единице. Поскольку каждое из них не
превышает единицы, любое слагаемое с положительным весом также
должно давать единицу. Равенство в неравенстве Коши достигается
только при совпадении направлений. Поэтому для каждого ненулевого
оптимального импульса
\begin {equation}
 \lVert\boldsymbol p(\varphi_j)\rVert=1,
 \qquad
 \Delta\widehat{\boldsymbol V}_j^{\text{\normalfont имп}}={\Delta\widehat V_j^{\text{\normalfont имп}}}\boldsymbol p(\varphi_j),
 \qquad {\Delta\widehat V_j^{\text{\normalfont имп}}}>0.
 \label {eq:primer-contact}
\end {equation}
Таким образом, импульсы располагаются в точках абсолютного максимума
модуля базис-вектора и направлены вдоль него. Все используемые
максимумы имеют одну высоту, равную единице.
Условие \eqref{eq:primer-contact} задаёт \emph{направление}, но не
величину импульса. Величины ${\Delta\widehat V_j^{\text{\normalfont имп}}}$ определяются из граничных уравнений.

Вывод через выпуклую оболочку учитывает произвольное число импульсов.
Это устраняет существенный пробел, который возник бы при
дифференцировании задачи с заранее заданным числом манёвров:
стационарность внутри такой схемы сама по себе не исключает её
улучшения добавлением нового импульса.

\subsection{Первая вариация по величинам, направлениям и положениям импульсов}

Условия в используемых точках можно дополнительно получить
непосредственным варьированием. Для фиксированного числа ненулевых
импульсов составим функцию Лагранжа
\begin {equation}
 \mathcal L=
 \sum_j\lVert\Delta\widehat{\boldsymbol V}_j^{\text{\normalfont имп}}\rVert
 -\boldsymbol l^{\mathsf T}
       \left(\sum_j B(\varphi_j)\Delta\widehat{\boldsymbol V}_j^{\text{\normalfont имп}}-\boldsymbol b\right).
\end {equation}
Поскольку
$\delta\lVert\Delta\widehat{\boldsymbol V}_j^{\text{\normalfont имп}}\rVert
  =\bigl(\Delta\widehat{\boldsymbol V}_j^{\text{\normalfont имп}}\bigr)^{\mathsf T}\delta\Delta\widehat{\boldsymbol V}_j^{\text{\normalfont имп}}/{\Delta\widehat V_j^{\text{\normalfont имп}}}$,
первая вариация имеет вид
\begin {equation}
 \delta\mathcal L=
 \sum_j\left(\frac{{\Delta\widehat{\boldsymbol V}_j^{\text{\normalfont имп}}}}{{{\Delta\widehat V_j^{\text{\normalfont имп}}}}}
             -\boldsymbol p(\varphi_j)\right)^{\mathsf T}
                \delta\Delta\widehat{\boldsymbol V}_j^{\text{\normalfont имп}}
 -\sum_j \boldsymbol p'(\varphi_j)^{\mathsf T}
                  \Delta\widehat{\boldsymbol V}_j^{\text{\normalfont имп}}\,\delta\varphi_j,
 \label {eq:impulse-first-variation}
\end {equation}
где штрих означает производную по $\varphi$.
Произвольность вариаций трёх компонент $\Delta\widehat{\boldsymbol V}_j^{\text{\normalfont имп}}$ даёт
\eqref{eq:primer-contact}; произвольность вариации внутренней фазы
даёт
\begin {equation}
 \boldsymbol p(\varphi_j)^{\mathsf T}\boldsymbol p'(\varphi_j)=0,
 \qquad
 \left.\frac{{\dd}}{{\dd\varphi}}
     \lVert\boldsymbol p(\varphi)\rVert\right|_{\varphi=\varphi_j}=0.
 \label {eq:primer-stationarity}
\end {equation}
Производная самого вектора при этом, вообще говоря, не равна нулю.
Базис-вектор может поворачиваться, сохраняя свой модуль.
Кроме того, стационарная точка может оказаться минимумом либо
не самым высоким максимумом. Поэтому
\eqref{eq:primer-stationarity} не заменяет глобальную проверку
\eqref{eq:primer-global-bound}.

Роль дополнительного импульса видна из его вклада в функцию Лагранжа.
При добавлении $\varepsilon\boldsymbol d$ в фазе $\varphi$, где
$\varepsilon>0$ и $\lVert\boldsymbol d\rVert=1$, вклад равен
\begin {equation}
 \varepsilon\bigl(1-
       \boldsymbol p(\varphi)^{\mathsf T}\boldsymbol d\bigr).
\end {equation}
Наиболее выгодное направление совпадает с $\boldsymbol p$; тогда
коэффициент при $\varepsilon$ равен
$1-\lVert\boldsymbol p(\varphi)\rVert$.
В данном выводе общий множитель $\boldsymbol l$ и невозможность
отрицательного коэффициента уже обоснованы
\eqref{eq:supporting-plane}. Одна только подстановка нового импульса
в исходное управление была бы недостаточна: она нарушает конечные
условия, которые также необходимо учитывать.

Если импульсы допускаются лишь на заданном временном интервале,
в его фиксированных концах двусторонняя вариация времени отсутствует.
Тогда равенство \eqref{eq:primer-stationarity} для концевого импульса
не обязательно: возникают односторонние условия.
В нашей задаче разрешена любая фаза на окружности, поэтому такого
конца допустимого интервала фаз нет.

\subsection{Достаточность и нижняя оценка характеристической скорости}

\begin{theorem}
Для линейной задачи \eqref{eq:linear-optimization} с
$\boldsymbol b\ne0$ набор импульсов глобально оптимален тогда и
только тогда, когда существуют постоянные $\boldsymbol l$,
для которых выполнены условия
\begin {equation}
 \boxed{\begin {gathered}
   \sum_jB(\varphi_j)\Delta\widehat{\boldsymbol V}_j^{\text{\normalfont имп}}=\boldsymbol b,\\
   \boldsymbol p(\varphi)=B(\varphi)^{\mathsf T}\boldsymbol l,
   \qquad \lVert\boldsymbol p(\varphi)\rVert\leq1
                              \quad\text {при всех }\varphi,\\
   \Delta\widehat{\boldsymbol V}_j^{\text{\normalfont имп}}={\Delta\widehat V_j^{\text{\normalfont имп}}}\boldsymbol p(\varphi_j),
   \qquad {\Delta\widehat V_j^{\text{\normalfont имп}}}>0,\qquad
   \lVert\boldsymbol p(\varphi_j)\rVert=1.
 \end {gathered}}
 \label {eq:certificate}
\end {equation}
\end{theorem}

\begin{proof}
Необходимость установлена построением опорной гиперплоскости.
Для доказательства достаточности рассмотрим любой другой
допустимый набор $\{\varphi_k^{\mathrm {alt}},\Delta\widehat{\boldsymbol V}_k^{\text{\normalfont имп},\,\mathrm{alt}}\}$,
в том числе с другим числом импульсов. Имеем
\begin {align}
 J^{\mathrm {alt}}
 &=\sum_k\lVert\Delta\widehat{\boldsymbol V}_k^{\text{\normalfont имп},\,\mathrm{alt}}\rVert
 \geq\sum_k\boldsymbol p(\varphi_k^{\mathrm {alt}})^{\mathsf T}
                      \Delta\widehat{\boldsymbol V}_k^{\text{\normalfont имп},\,\mathrm{alt}} \notag\\
 &=\boldsymbol l^{\mathsf T}
         \sum_k B(\varphi_k^{\mathrm {alt}})\Delta\widehat{\boldsymbol V}_k^{\text{\normalfont имп},\,\mathrm{alt}}
  =\boldsymbol l^{\mathsf T}\boldsymbol b.
 \label {eq:dual-lower-bound}
\end {align}
Для проверяемого набора все неравенства обращаются в равенства:
\begin {equation}
 \boldsymbol l^{\mathsf T}\boldsymbol b
   =\sum_j{\Delta\widehat V_j^{\text{\normalfont имп}}}\lVert\boldsymbol p(\varphi_j)\rVert^2
   =\sum_j{\Delta\widehat V_j^{\text{\normalfont имп}}}=J.
\end {equation}
Следовательно, $J^{\mathrm {alt}}\geq J$ для любого допустимого управления.
\end{proof}

Это доказательство использует линейность граничных уравнений и
выпуклость нормы. Для точного нелинейного перехода те же условия
на базис-вектор в общем случае являются лишь необходимыми:
линейная оценка \eqref{eq:dual-lower-bound} уже не связывает
произвольные конечные орбиты с одним и тем же оператором $B$.

Полезна также двойственная запись оптимальной цены:
\begin {equation}
 J_*(\boldsymbol b)=
 \max_{\boldsymbol l}
 \left\{\boldsymbol l^{\mathsf T}\boldsymbol b:
       \max_\varphi\lVert B(\varphi)^{\mathsf T}\boldsymbol l\rVert
           \leq1\right\}.
 \label {eq:dual-problem}
\end {equation}
Если $J_*$ дифференцируема по всем пяти компонентам
$\boldsymbol b$, то оптимальный множитель равен её градиенту:
\begin {equation}
 \boldsymbol l=\nabla_{\boldsymbol b}J_*.
 \label {eq:cost-gradient}
\end {equation}
Действительно, тот же допустимый множитель даёт для изменённого
задания нижнюю оценку
$J_*(\boldsymbol b+\delta\boldsymbol b)
 \geq J_*(\boldsymbol b)+\boldsymbol l^{\mathsf T}\delta\boldsymbol b$.
В точке дифференцируемости такая линейная опорная функция имеет
коэффициенты градиента. На границе ветвей множитель может быть
неединственным; тогда использовать производную несуществующей
гладкой ветви нельзя.

Нулевое значение четвёртой компоненты задания не означает
$l_s=0$. Производная в \eqref{eq:cost-gradient} берётся по
полному вектору изменений, включая обе компоненты наклонения.
Если до дифференцирования повернуть оси вдоль линии узлов,
при вариации направления конечной плоскости необходимо учесть
и изменение этого поворота.

\subsection{Связь с сопряжёнными уравнениями Лоудена}

Покажем, откуда появляется та же векторная функция в общей задаче.
В этом подразделе $\boldsymbol R$ и $\boldsymbol V$ --- размерные
инерциальные векторы, $R=\lVert\boldsymbol R\rVert$, а $\Delta\boldsymbol V_j^{\text{\normalfont имп}}$ --- размерные
импульсы. На баллистических дугах
\begin {equation}
 \dot{\boldsymbol R}=\boldsymbol V,
 \qquad \dot{\boldsymbol V}=\boldsymbol g(\boldsymbol R),
 \qquad
 \boldsymbol g(\boldsymbol R)=-\frac{{\mu}}{{R^3}}\boldsymbol R.
\end {equation}
Минимизируется $\sum_j\lVert\Delta\boldsymbol V_j^{\text{\normalfont имп}}\rVert$.
Для регулярной, или нормальной, экстремали коэффициент при этом
функционале в расширенной вариационной задаче можно принять равным
единице. Нормальность означает, что множители не описывают только
вырожденные ограничения при нулевом коэффициенте стоимости.
Для линейной задачи необходимость условий уже доказана выше
без дополнительного предположения о нормальности.

Введём сопряжённые векторы $\boldsymbol a(t)$ и
$\boldsymbol p(t)$ и добавим к стоимости интегралы по всем
баллистическим дугам:
\begin {equation}
 \mathcal A=\sum_j\lVert\Delta\boldsymbol V_j^{\text{\normalfont имп}}\rVert
 +\sum_{\text{дуги}}\int
 \left[
 \boldsymbol a^{\mathsf T}(\dot{\boldsymbol R}-\boldsymbol V)
 +\boldsymbol p^{\mathsf T}
          (\dot{\boldsymbol V}-\boldsymbol g(\boldsymbol R))
 \right]\dd t.
 \label {eq:augmented-coast-action}
\end {equation}
Множители концевых ограничений здесь не выписаны, поскольку сначала
рассматриваются внутренние вариации. Варьирование $\boldsymbol R$
и $\boldsymbol V$ и интегрирование слагаемых с производными по
частям дают внутри дуг
\begin {equation}
 \dot{\boldsymbol a}
       =-\left(\frac{{\partial\boldsymbol g}}{{\partial\boldsymbol R}}\right)^{\mathsf T}
          \boldsymbol p,
 \qquad
 \dot{\boldsymbol p}=-\boldsymbol a.
\end {equation}
Исключая $\boldsymbol a$, получаем уравнение базис-вектора:
\begin {equation}
 \ddot{\boldsymbol p}
  =\left(\frac{{\partial\boldsymbol g}}{{\partial\boldsymbol R}}\right)^{\mathsf T}
       \boldsymbol p
  =-\frac{{\mu}}{{R^3}}\boldsymbol p
    +\frac{{3\mu}}{{R^5}}
           (\boldsymbol p^{\mathsf T}\boldsymbol R)\boldsymbol R.
 \label {eq:lawden-primer-ode}
\end {equation}
Второе равенство следует из покомпонентного дифференцирования
$g_i=-\mu R_i/R^3$:
$\partial g_i/\partial R_k
 =-\mu\delta_{ik}/R^3+3\mu R_i R_k/R^5$.
Матрица производных симметрична.
Уравнение \eqref{eq:lawden-primer-ode} линейно по
$\boldsymbol p$ вдоль заданной баллистической дуги. Сама задача
поиска этой дуги между двумя орбитами от этого линейной не становится.

В момент импульса
$\boldsymbol R^+=\boldsymbol R^-$ и
$\boldsymbol V^+=\boldsymbol V^-+\Delta\boldsymbol V^{\text{\normalfont имп}}$.
Вариации положения и доимпульсной скорости в граничных членах
\eqref{eq:augmented-coast-action} входят в выражение
\begin {equation}
 \begin {aligned}
 \delta\mathcal A_j={}&(\boldsymbol a^--\boldsymbol a^+)^{\mathsf T}
 \delta\boldsymbol R
 +(\boldsymbol p^--\boldsymbol p^+)^{\mathsf T}\delta\boldsymbol V^-\\
 &+\left(\frac{{\Delta\boldsymbol V^{\text{\normalfont имп}}}}{{\lVert\Delta\boldsymbol V^{\text{\normalfont имп}}\rVert}}
 -\boldsymbol p^+\right)^{\mathsf T}\delta\Delta\boldsymbol V^{\text{\normalfont имп}}.
 \end {aligned}
\end {equation}
Здесь использовано равенство
$\delta\boldsymbol V^+=\delta\boldsymbol V^-+\delta\Delta\boldsymbol V^{\text{\normalfont имп}}$.
Независимость трёх указанных вариаций даёт
$\boldsymbol a^+=\boldsymbol a^-$ и
$\boldsymbol p^+=\boldsymbol p^-$.
Коэффициент при вариации самого импульса равен
\begin {equation}
 \frac{{\Delta\boldsymbol V^{\text{\normalfont имп}}}}{{\lVert\Delta\boldsymbol V^{\text{\normalfont имп}}\rVert}}
          -\boldsymbol p.
\end {equation}
Следовательно, $\boldsymbol p$ направлен вдоль ненулевого импульса
и имеет в этой точке единичный модуль. Поскольку
$\dot{\boldsymbol p}=-\boldsymbol a$, первая инерциальная производная
базис-вектора также непрерывна. Утверждать непрерывность всех его
производных оснований нет.

Для свободного внутреннего момента импульса возникает условие
непрерывности гамильтониана
\begin {equation}
 H=\boldsymbol a^{\mathsf T}\boldsymbol V
            +\boldsymbol p^{\mathsf T}\boldsymbol g.
\end {equation}
Оно следует из коэффициента при перемещении границы двух
баллистических интегралов. Так как положение и сопряжённые векторы
непрерывны,
\begin {equation}
 0=H^+-H^-
  =\boldsymbol a^{\mathsf T}\Delta\boldsymbol V^{\text{\normalfont имп}}
  =-\dot{\boldsymbol p}^{\mathsf T}\Delta\boldsymbol V^{\text{\normalfont имп}}.
\end {equation}
После подстановки направления импульса получаем
$\boldsymbol p^{\mathsf T}\dot{\boldsymbol p}=0$, то есть
стационарность его модуля. Это размерный инерциальный аналог
\eqref{eq:primer-stationarity}.

Наконец, для малого дополнительного импульса $\Delta\boldsymbol V^{\text{\normalfont имп}}$
точная первая вариация расширенного функционала для нормальной
экстремали содержит
$\lVert\Delta\boldsymbol V^{\text{\normalfont имп}}\rVert-
\boldsymbol p^{\mathsf T}\Delta\boldsymbol V^{\text{\normalfont имп}}$.
Её неотрицательность при произвольном направлении даёт
$\lVert\boldsymbol p\rVert\leq1$.
Таким образом, условия Лоудена согласуются с полученными выше
условиями, но их достаточность для нелинейной задачи из этой первой
вариации не следует.

\subsection{Базис-вектор на опорной круговой орбите}

Положим теперь $R=r_0$ и рассмотрим вращающиеся единичные векторы
$\boldsymbol e_r$, $\boldsymbol e_t$, $\boldsymbol e_z$,
для которых
\begin {equation}
 \dot{\boldsymbol e}_r=n_0\boldsymbol e_t,
 \qquad \dot{\boldsymbol e}_t=-n_0\boldsymbol e_r,
 \qquad \dot{\boldsymbol e}_z=0.
\end {equation}
Компоненты $p_r,p_t,p_z$ определены разложением
$\boldsymbol p=p_r\boldsymbol e_r+p_t\boldsymbol e_t+p_z\boldsymbol e_z$.
При дифференцировании необходимо учитывать вращение осей:
\begin {equation}
 \dot{\boldsymbol p}
 = (\dot p_r-n_0p_t)\boldsymbol e_r
  +(\dot p_t+n_0p_r)\boldsymbol e_t
  +\dot p_z\boldsymbol e_z.
\end {equation}
Подстановка второй производной в
\eqref{eq:lawden-primer-ode} и переход от времени к
$\varphi=n_0t+\varphi_0$ дают
\begin {equation}
 p_r''-2p_t'-3p_r=0,
 \qquad p_t''+2p_r'=0,
 \qquad p_z''+p_z=0.
 \label {eq:circular-primer-ode}
\end {equation}
Из второго уравнения $p_t'+2p_r=C$.
Тогда первое превращается в $p_r''+p_r=2C$.
Общее решение имеет шесть постоянных:
\begin {equation}
 \begin {aligned}
 p_r&=c_1\cos\varphi+c_2\sin\varphi+2C,\\
 p_t&=2c_2\cos\varphi-2c_1\sin\varphi-3C\varphi+c_3,\\
 p_z&=c_4\cos\varphi+c_5\sin\varphi.
 \end {aligned}
 \label {eq:circular-primer-general}
\end {equation}

Появление секулярного члена $-3C\varphi$ не противоречит
\eqref{eq:primer-components}: ещё не использовано условие
трансверсальности по свободному времени.
В автономной задаче с конечным множеством, не зависящим явно от
времени, и свободным конечным временем это условие есть $H=0$.
Поясним появление этого равенства. Пусть
$\boldsymbol Z=(\boldsymbol R,\boldsymbol V)$ и
$\boldsymbol \Lambda=(\boldsymbol a,\boldsymbol p)$.
Полная вариация конечного состояния связана с вариацией при
фиксированном времени равенством
$\delta\boldsymbol Z_f=\delta\boldsymbol Z(t_f)+
\dot{\boldsymbol Z}(t_f)\delta t_f$.
Поэтому конечный член после интегрирования по частям имеет вид
\[
 \boldsymbol \Lambda_f^{\mathsf T}\delta\boldsymbol Z_f-H_f\delta t_f.
\]
Трансверсальность по конечному состоянию означает, что
$\boldsymbol \Lambda_f$ нормален к допустимому конечному многообразию.
Его скалярное произведение с любой касательной вариацией
$\delta\boldsymbol Z_f$ равно нулю. Так как само многообразие
от времени не зависит, $\delta t_f$ независимо и произвольно:
остаётся $H_f=0$. На пассивных дугах автономной задачи $H$ постоянно,
а при свободных внутренних импульсах непрерывно; следовательно, $H=0$
на всей экстремали. В нашей задаче конечное многообразие является
заданной орбитой со свободной фазой. Её инвариантность относительно
пассивного движения также согласуется с условием $H=0$.
На круговой орбите $\boldsymbol V=V_0\boldsymbol e_t$,
$\boldsymbol g=-n_0V_0\boldsymbol e_r$ и
$\boldsymbol a=-\dot{\boldsymbol p}$, поэтому
\begin {equation}
 H=-n_0V_0(p_t'+2p_r)=-n_0V_0 C.
\end {equation}
Следовательно, $C=0$. Сравнение
\eqref{eq:circular-primer-general} с
\eqref{eq:primer-components} даёт
\begin {equation}
 c_1=-l_y,\qquad c_2=l_x,\qquad c_3=2l_a,
 \qquad c_4=l_c,\qquad c_5=l_s.
\end {equation}
Итак, пятиконстантная форма базис-вектора следует и из точных
сопряжённых уравнений после перехода к опорному кругу и использования
свободного времени. Круговая орбита сама по себе не требует $C=0$.
Для встречи с заданной фазой в заданное время условия
трансверсальности и число граничных ограничений следует выводить
заново.

\subsection{Геометрия годографа и типы оптимальных переходов}
\label{subsec:primer-geometry}

Вначале предположим
$R=\sqrt{l_x^2+l_y^2}>0$ и положим
\begin {equation}
 \begin {gathered}
 \alpha=\operatorname {atan2}(l_y,l_x),
 \qquad \psi=\varphi-\alpha,\qquad D=2l_a,\\
 E=l_s\sin\alpha+l_c\cos\alpha,
 \qquad F=l_s\cos\alpha-l_c\sin\alpha.
 \end {gathered}
\end {equation}
Тогда
\begin {equation}
 \boldsymbol p(\psi)=
 \begin {pmatrix}R\sin\psi\\D+2R\cos\psi\\E\cos\psi+F\sin\psi\end {pmatrix}
 =\begin {pmatrix}0\\D\\0\end {pmatrix}
 +\begin {pmatrix}0\\2R\\E\end {pmatrix}\cos\psi
 +\begin {pmatrix}R\\0\\F\end {pmatrix}\sin\psi.
 \label {eq:rotated-primer}
\end {equation}
Годограф в пространстве трёх \emph{вращающихся компонент} является
эллипсом с центром $(0,D,0)$. Два вектора при синусе и косинусе
задают его плоскость. Его радиально-трансверсальная проекция имеет
полуоси $R$ и $2R$.
Это вспомогательное пространство компонент, а не физическая
траектория конца вектора в неподвижных координатах.

Квадрат модуля равен
\begin {equation}
 \begin {aligned}
 f(\psi)=\lVert\boldsymbol p(\psi)\rVert^2
   &=D^2+R^2+F^2
    +(3R^2+E^2-F^2)\cos^2\psi\\
   &\quad+2EF\sin\psi\cos\psi+4DR\cos\psi.
 \end {aligned}
 \label {eq:primer-norm-quadratic}
\end {equation}
Импульсы возможны только в точках касания этого эллипса с единичной
сферой, причём весь эллипс должен лежать внутри сферы.
Положения и число точек касания определяют классификацию,
используемую Эдельбаумом и Барановым
\cite{edelbaum,baranov-ru,baranov-en}.

\paragraph{Центральный эллипс: узловой случай (NC).}
Если $D=0$, то
$\boldsymbol p(\psi+\pi)=-\boldsymbol p(\psi)$.
У некругового эллипса наиболее удалены от начала координат концы
его большой оси. Таким образом, возможны два противоположных
максимума; их фазы отличаются на $\pi$, а все локальные компоненты
направлений импульсов противоположны.
Для вычисления максимума достаточно максимального собственного
значения симметричной матрицы
\begin {equation}
 Q=\begin {pmatrix}
       4R^2+E^2&EF\\ EF&R^2+F^2
    \end {pmatrix},
 \qquad
 f=(\cos\psi,\sin\psi)Q
        \begin {pmatrix}\cos\psi\\\sin\psi\end {pmatrix}.
\end {equation}
Если оно равно единице, эллипс касается сферы концами большой оси.
Название ``узловой'' будет обосновано при подстановке граничных
условий пространственного перехода: для ${\Delta i}>0$ используемые точки
окажутся на линии взаимных узлов.

\paragraph{Смещённый эллипс: невырожденный случай (NDC).}
Пусть $D\ne0$ и имеются две различные точки абсолютного максимума.
Покажем, почему они должны быть симметричны относительно $\psi=0$.
Функция $1-f(\psi)$ --- неотрицательный тригонометрический многочлен
степени два. В каждой точке касания со сферой он имеет двойной нуль.
Если два нуля обозначить $m-\delta$ и $m+\delta$, такой многочлен
представляется в виде
\begin {equation}
 1-f(\psi)=K[\cos(\psi-m)-\cos\delta]^2,
 \qquad K>0.
 \label {eq:contact-factorization}
\end {equation}
Действительно, произведение
$\sin^2((\psi-m+\delta)/2)
 \sin^2((\psi-m-\delta)/2)$ имеет именно эти два двойных нуля;
раскрытие произведения даёт правую часть с постоянным множителем.
Многочлен степени два с такими нулями определяется с точностью до
этого множителя.
Первая гармоника в \eqref{eq:primer-norm-quadratic} содержит только
$\cos\psi$ и имеет ненулевой коэффициент $4DR$.
Сравнение с \eqref{eq:contact-factorization} даёт
$m=0$ по модулю $\pi$. Следовательно, точки можно записать как
$\psi=\pm\psi_*$. Коэффициент при $\sin2\psi$ должен обращаться
в нуль, то есть $EF=0$.

Если бы $F=0$, зависимость от $c=\cos\psi$ была бы выпуклой
квадратичной функцией с коэффициентом $3R^2+E^2>0$.
При $DR\ne0$ она имеет единственный абсолютный максимум на
$[-1,1]$, что противоречит наличию двух различных точек касания.
Значит, $E=0$, и два максимума существуют при
\begin {equation}
 F^2>3R^2,\qquad
 c_*:=\cos\psi_*=\frac{{2DR}}{{F^2-3R^2}},
 \qquad |c_*|<1.
 \label {eq:ndc-primer-angles}
\end {equation}
Эти формулы получаются дифференцированием
$f(c)=D^2+R^2+F^2-(F^2-3R^2)c^2+4DRc$.
Условие нормировки имеет вид
\begin {equation}
 D^2+R^2+F^2+\frac{{4D^2R^2}}{{F^2-3R^2}}=1.
 \label {eq:ndc-primer-normalization}
\end {equation}
В двух точках касания трансверсальные компоненты базис-вектора
одинаковы, а радиальные и нормальные компоненты противоположны.
При $|c_*|=1$ точки сливаются; это граничный случай с возможным
единственным ненулевым импульсом.

\paragraph{Окружность: сингулярный случай (SC).}
Если центральный эллипс имеет постоянный единичный радиус,
разделение на отдельные точки максимума исчезает.
Из $D=0$ и $Q$ равной единичной матрице получаем
\begin {equation}
 \begin {gathered}
 E=0,\qquad R=\frac{1}{2},\qquad F=\pm\frac{{\sqrt3}}{{2}},\\
 \boldsymbol p(\psi)=
 \begin {pmatrix}
   \tfrac{1}{2}\sin\psi\\ \cos\psi\\
   \pm\tfrac{{\sqrt3}}{{2}}\sin\psi
 \end {pmatrix},
 \qquad \lVert\boldsymbol p(\psi)\rVert\equiv1.
 \end {gathered}
 \label {eq:singular-primer-circle}
\end {equation}
Все фазы имеют одинаковую предельную эффективность, но направления
импульсов по-прежнему заданы функцией $\boldsymbol p(\psi)$.
Свобода выбора фаз не снимает граничных условий и требования
неотрицательности величин ${\Delta\widehat V_j^{\text{\normalfont имп}}}$.
Различные допустимые распределения импульсов, удовлетворяющие
\eqref{eq:certificate}, имеют одинаковую оптимальную стоимость.
В частности, постоянным здесь является модуль, а не сам вектор.

\paragraph{Вырождение в отрезок и точку.}
При $R=0$ угол $\alpha$ не определён; делить на $R$ нельзя.
В исходных координатах
\begin {equation}
 \boldsymbol p(\varphi)=
      (0,D,l_s\sin\varphi+l_c\cos\varphi)^{\mathsf T}.
 \label {eq:line-primer-general}
\end {equation}
Если $l_s^2+l_c^2>0$, годограф --- отрезок, и
$\max\lVert\boldsymbol p\rVert^2=D^2+l_s^2+l_c^2$.
Два конца отрезка соответствуют противоположным фазам;
трансверсальные компоненты одинаковы, нормальные противоположны,
радиальные равны нулю.
При $l_s=l_c=0$ годограф вырождается в точку, а нормировка даёт
$D=\pm1$. Этот предел особенно существенен для компланарных
переходов с преобладающим изменением большой полуоси.

Таким образом, геометрия базис-вектора определяет допустимые типы
расположения импульсов. Она ещё не определяет, какой тип отвечает
заданным ${\Delta a},{\Delta e_x},{\Delta e_y},{\Delta i}$: это устанавливается совместным решением
граничных уравнений, условий касания и неравенств на величины
импульсов. Именно эта задача решается в следующем разделе.

% ==================== planar ====================
